\documentclass[11pt]{article}
\usepackage[a4paper,margin=25mm]{geometry}
\usepackage[T1]{fontenc}
\usepackage{lmodern,microtype}
\usepackage{amsmath,amssymb,amsthm,mathtools}
\usepackage{booktabs,array}
\usepackage[hidelinks]{hyperref}
\usepackage{fancyhdr}
\hypersetup{pdftitle={Occupied-fibre rigidity and entropy tail nonattainment},pdfauthor={B2 research continuation},pdfsubject={Continuation III for Q4296 and Q4660}}
\setlength{\parindent}{0pt}
\setlength{\parskip}{3pt}
\setlength{\emergencystretch}{2em}
\setlength{\headheight}{14pt}
\pagestyle{fancy}\fancyhf{}
\fancyhead[L]{\small B2 research continuation III}
\fancyhead[R]{\small 13 October 2026}
\fancyfoot[C]{\thepage}
\newtheorem{theorem}{Theorem}[section]
\newtheorem{lemma}[theorem]{Lemma}
\newtheorem{proposition}[theorem]{Proposition}
\newtheorem{corollary}[theorem]{Corollary}
\theoremstyle{remark}\newtheorem{remark}[theorem]{Remark}
\newcommand{\F}{\mathbb F_3}
\newcommand{\N}{\mathbb N_0}
\newcommand{\R}{\mathbb R}
\newcommand{\Prob}{\mathbb P}
\newcommand{\E}{\mathbb E}
\newcommand{\supp}{\operatorname{supp}}
\newcommand{\Dplus}{\Delta_+}
\newcommand{\Dminus}{\Delta_-}
\title{\vspace{-12mm}\textbf{Occupied-fibre rigidity and\ entropy tail nonattainment}\\[4pt]\large New structural results for the B2 midpoint and entropy programs}
\author{Research continuation from the supplied handoff}
\date{13 October 2026}
\begin{document}
\maketitle\vspace{-6mm}
\begin{abstract}
An admissible two-colour remainder in $\F^n$ containing no affine plane has
at most $\frac52\,2^n-4$ points for $n\geq3$. The three-dimensional bound of
$16$ is sharp. The proof establishes a normal form relative to a
maximum-dimensional remainder flat: all but at most one nonfull fibre
have monochromatic affine coloured complements, and the set of occupied
fibres inherits an admissible two-colouring. For the sum--difference entropy
ratio, a full-support law on $\N$ with
$-\log p_n=c n^\beta+o(n^\beta)$ can be stationary at ratio $\rho>1$ only if
$\beta=\log_2(4\rho/(1+\rho))$. Every prescribed positive tail exponent and
coefficient nevertheless give a class with the same supremum as the
unrestricted problem. In particular the one-sided Gaussian-tail class has
the full sharp supremum but no maximizer. No new numerical entropy lower
bound or general many-colour codimension improvement is claimed.
\end{abstract}

\section{Scope and baselines}

A partition $\F^n=A\sqcup B\sqcup R$ is \emph{admissible} if
\[
 x,y\in A\ \text{or}\ x,y\in B \quad\Longrightarrow\quad (x+y)/2\notin R.
\]
An affine plane has dimension two and nine points. A cap contains no affine
line. Empty colour classes are allowed. Density assumptions will be stated
separately rather than built into this definition.

The targets are Q4296 (midpoint colourings) and Q4660 (entropy ratio);
their stable identifiers are retained in the result ledger. The primary
midpoint paper gives an
$O(k^2)$ codimension guarantee at density $1/10$ \cite{midpoint}.
Continuation I \cite{first} establishes $O(k\log(2k))$ using polynomial
Freiman--Ruzsa covering. That general bound is unchanged here.
Continuation II \cite{previous} establishes the two-colour cap classification
and the entropy lower bound
\begin{equation}\label{eq:knownrange}
 1.4051810329<C_*\leq2.
\end{equation}
The upper bound is the established entropic Ruzsa triangle inequality;
its sharpness is the question in Section 8 of \cite{li}. The catalog's
reference to Conjecture 1.10 concerns a different, sum--product statement.
That provenance correction is retained separately without modifying the
original records.

The present results have complete written proofs. Exact finite tests are
controls on their elementary components, not formal proofs of the
all-dimension or infinite-support theorems. Worldwide priority, external
peer review, and proof-assistant verification have not been established.
Neither complete catalog problem is declared solved.

\section{Normal form along a maximum-dimensional remainder flat}

We use the following result of Continuation II. A short proof is reproduced
in Appendix~\ref{app:cap} to make the present midpoint arguments self-contained.

\begin{lemma}[Prior cap classification]\label{lem:cap}
A cap that is an admissible two-colour remainder in $\F^n$ is empty, a
singleton, two points, or a four-point affine parallelogram. In particular,
its size is at most four. Here an affine parallelogram is
$a+\{0,u,v,u+v\}$ with $u,v$ linearly independent.
\end{lemma}

Suppose $R\ne\varnothing$, and let $d$ be the largest dimension of an affine
subspace contained in $R$. Fix the direction $V\leq\F^n$ of one such
$d$-dimensional subspace. Let $\pi:\F^n\to Q=\F^n/V$, and write
\[
 F_q=\pi^{-1}(q),\qquad R_q=R\cap F_q,\qquad C_q=F_q\setminus R.
\]
The full and occupied quotient sets are respectively
\[
 S=\{q:F_q\subseteq R\},\qquad D=\{q:R_q\ne\varnothing\}.
\]
Thus $S\subseteq D$, and $S$ is nonempty.

\begin{theorem}[Maximal-flat fibre rigidity]\label{thm:normal}
In the notation above, $S$ has size $1$, $2$, or $4$, and in the last case
is an affine parallelogram. All but at most one fibre outside $S$ have a
nonempty monochromatic coloured part $C_q$, and that part is an affine
subspace of $F_q$. Consequently, in every such regular fibre,
\begin{equation}\label{eq:fibresize}
 |R_q|=3^d-3^j \quad\text{for some }0\leq j\leq d.
\end{equation}
If $|S|=1$ there is no exceptional fibre. If $|S|=2$ or $4$, the only
possible exception is the fibre over the unique centre of symmetry of $S$.
\end{theorem}
\begin{proof}
First, $S$ inherits an admissible colouring: for every $q\notin S$, choose a
representative $x_q\in F_q\setminus R$ and assign its colour to $q$. A
same-colour pair whose quotient midpoint lies in $S$ would have its actual
midpoint in a fibre entirely contained in $R$, a contradiction. A line in
$S$ would lift to a $(d+1)$-dimensional affine subspace in $R$. Thus $S$ is
a cap and Lemma~\ref{lem:cap} applies.

On $Q\setminus S$, join distinct $q,t$ when $(q+t)/2\in S$. In that case
every point of $C_q$ must have colour different from every point of $C_t$.
Both sets are nonempty. With only two available colours, both are therefore
monochromatic.

A vertex $q\notin S$ is isolated precisely when $2s-q\in S$ for all $s\in S$.
In characteristic three, $2s-q=2q-s$, so this says
\[
 \sigma_q(S)=S,\qquad \sigma_q(s)=2q-s.
\]
Two distinct symmetry centres would give a nonzero translation period of
$S$, by composing the reflections. Its three-point orbits would be affine
lines in $S$, contrary to the cap property. There is at most one isolated
vertex. A singleton's centre is its own point, not an exterior vertex;
a two-point set or parallelogram has exactly one centre, lying outside it.

Every non-isolated $C_q$ is monochromatic. The midpoint of two of its points
is in the same fibre and cannot be in $R$, so $C_q$ is midpoint-closed. A
nonempty midpoint-closed subset of a vector space over $\F$ is affine:
translate a member to zero, use midpoint$(0,x)=-x$ to obtain negatives, and
then midpoint$(x,y)=-x-y$ to obtain addition. This proves
\eqref{eq:fibresize} and the theorem.
\end{proof}

\begin{theorem}[Occupied-fibre inheritance]\label{thm:occupied}
The occupied set $D$ is itself the remainder of an admissible two-colour
partition of $Q$.
\end{theorem}
\begin{proof}
For $q\notin D$ the entire fibre $F_q$ is coloured. If it is regular,
Theorem~\ref{thm:normal} says that it is monochromatic; use that colour for
$q$. In the possible exceptional fibre use any colour that occurs there.

Suppose distinct $q,t\notin D$ receive the same colour $i$. At most one is
exceptional, so one of their entire fibres has colour $i$ and the other has
at least one point of colour $i$. Midpoints of that point with the entire
first fibre fill $F_{(q+t)/2}$. Hence the whole midpoint fibre avoids $R$,
and $(q+t)/2\notin D$. When $q=t$ the conclusion is automatic. This is the
required admissibility.
\end{proof}

\begin{remark}
The occupied-fibre conclusion is stronger than inheritance for full fibres.
It uses both the two-colour hypothesis and a maximum-dimensional remainder
direction. No assertion of occupied-fibre inheritance for arbitrary
quotients or arbitrary numbers of colours is made.
\end{remark}

\section{The sharp plane-free bound in three dimensions}

We record two elementary ingredients, distinguishing the classical one from
the present deductions.

\begin{lemma}[Classical affine blocking bound]\label{lem:block}
A set $C\subseteq\F^m$ meeting every affine hyperplane has at least $2m+1$
points. In particular, six points in $\F^3$ miss some affine plane.
\end{lemma}
\begin{proof}
This is the $q=3$ case of the Jamison--Brouwer--Schrijver blocking theorem
\cite{jamison,brouwer}; the polynomial proof is included for completeness.
Translate a point of $C$ to zero and define, over $\F$,
\[
 P(u)=\prod_{c\in C\setminus\{0\}}(1-u\cdot c).
\]
For $u\ne0$, the hyperplane $u\cdot x=1$ contains a point of
$C\setminus\{0\}$, so $P(u)=0$, whereas $P(0)=1$. Reduce $P$ modulo
$u_i^3-u_i$ in each variable. Reduced polynomials of degree at most two
in each variable represent functions on $\F^m$ uniquely, as follows by
successive one-variable interpolation. The reduced polynomial must be
$\prod_i(1-u_i^2)$, the indicator of zero. It has total degree $2m$;
reduction cannot increase total degree. Hence $|C|-1\geq2m$.
\end{proof}

\begin{lemma}[Small spanning sumsets in $\F^3$]\label{lem:sumset}
Four affinely independent points have sumset size ten. Five points spanning
$\F^3$ have sumset size at least eleven. A four-point planar cap has sumset
size nine, and an affine line together with an off-line point has sumset
size seven.
\end{lemma}
\begin{proof}
An affine change of coordinates puts an affine basis at
$E=\{0,e_1,e_2,e_3\}$. Its ten sums are
\[
 T=\{0,e_i,2e_i:1\leq i\leq3\}\ \cup\ \{e_i+e_j:i<j\}.
\]
For an additional point $v\notin E$, either $2v\notin T$, already adding a
new sum, or $v=-e_i$ or $v=-e_i-e_j$. In the first case $v+e_j\notin T$
for $j\ne i$, and in the second $v+e_k\notin T$ for the remaining index $k$.
This proves the five-point assertion. A four-point planar cap is a
parallelogram (Appendix~\ref{app:cap}); its sumset fills its plane. Finally,
for a line $L$ and $v\notin L$, the three parts $L+L$, $v+L$, and $\{2v\}$
are disjoint, of sizes three, three, and one.
\end{proof}

\begin{theorem}[Sharp three-dimensional threshold]\label{thm:three}
In an admissible two-colour partition of $\F^3$, a remainder of size at
least seventeen contains an affine plane. There is a plane-free remainder
of size sixteen.
\end{theorem}
\begin{proof}
Assume $|R|\geq17$, and put $C=A\sqcup B$. Then $|C|\leq10$. Relabel so
$a=|A|\geq|B|$. Admissibility gives
\[
 (A+A)/2\subseteq C.
\]
If $a\leq3$, then $|C|\leq6$, and Lemma~\ref{lem:block} gives a plane in $R$.

If $a=4$, then $|C|\leq8$. The affine-basis and planar-cap possibilities for
$A$ are excluded by Lemma~\ref{lem:sumset}, since their sumsets have sizes
ten and nine. Thus $A$ is a line together with an off-line point. Its
midpoint set gives seven points of $C$ in a plane $H$. At most one point
of $C$ lies outside $H$, so at least one of the other two planes parallel
to $H$ lies wholly in $R$.

If $a\geq5$ and $A$ spans three dimensions, it contains five spanning
points. Their sumset has size at least eleven, contradicting $|C|\leq10$.
Otherwise $A$ lies in a plane $H$ and has size at least five. Any two subsets
of $H$ with more than half its points intersect, so $(A+A)/2=H$.
Again $C$ has at most one point outside $H$, yielding a parallel plane in $R$.

For sharpness take
\[
 R=(\F\setminus\{0\})\times(\F^2\setminus\{0\}).
\]
Colour points with first coordinate zero by $A$, and the remaining points
with their last two coordinates zero by $B$. The partition is admissible,
and $|R|=2\cdot8=16$. If $L\subseteq R$ is affine, each of its two block
projections is an affine subspace avoiding zero. Such a projection lies in
a level set $\ell=1$ of a nonzero linear functional: write it as $a+W$
with $a\notin W$ and choose $\ell(W)=0$, $\ell(a)=1$. The two pulled-back
functionals have disjoint coordinate supports and are independent.
Consequently $\dim L\leq3-2=1$. A line is obtained by fixing the first
coordinate and one coordinate of the second block to one.
\end{proof}

\section{An all-dimension plane-free bound}

\begin{theorem}\label{thm:global}
For every $n\geq3$, an admissible two-colour remainder $R\subseteq\F^n$
containing no affine plane satisfies
\begin{equation}\label{eq:main}
 \boxed{\ |R|\leq20\,2^{n-3}-4=\frac52\,2^n-4.\ }
\end{equation}
\end{theorem}
\begin{proof}
Let $M_n$ be the maximum size of a plane-free admissible two-colour
remainder. Theorem~\ref{thm:three} gives $M_3=16$.

If $R$ has no line, Lemma~\ref{lem:cap} gives $|R|\leq4$. Otherwise a line
in $R$ is a maximum-dimensional remainder flat, since $R$ has no plane.
Choose its direction $V$. In the preceding notation $d=1$. Every regular
nonfull fibre has either zero or two remainder points, by
\eqref{eq:fibresize}, and full fibres have three. At most one occupied
fibre can have just one remainder point.

The occupied quotient $D\subseteq\F^{n-1}$ is admissible by
Theorem~\ref{thm:occupied}. It is also plane-free. Indeed, if a quotient
plane $U$ lay in $D$, its preimage would be an affine three-space containing
nine occupied line fibres. At least eight would contain at least two
remainder points and the possible exceptional fibre at least one. The
three-space would therefore contain at least seventeen remainder points,
contradicting Theorem~\ref{thm:three}.

Every nonfull fibre has at most two remainder points, and there are at most
four full fibres. Hence, for $n\geq4$,
\[
 |R|\leq2|D|+|S|\leq2M_{n-1}+4.
\]
The cap case is smaller as well. Iterating the recurrence from $M_3=16$
yields $M_n+4\leq2^{n-3}(16+4)$, proving \eqref{eq:main}.
\end{proof}

\begin{corollary}[Explicit density consequence]\label{cor:density}
An admissible two-colour remainder with $|R|\geq3^n/10$ contains an affine
plane whenever $n\geq8$. More generally a plane is forced when
\[
 |R|/3^n>\frac52\left(\frac23\right)^n-\frac4{3^n}.
\]
\end{corollary}
\begin{proof}
At $n=8$, the upper bound in \eqref{eq:main} is $636$, below
$3^8/10=656.1$. The bound divided by $3^n$ decreases for $n\geq3$:
its successive ratio is $2/3+4/(3B_n)<1$, where $B_n=20\,2^{n-3}-4\geq16$.
\end{proof}

\begin{remark}
The theorem gives a quantitative extension from line-free to plane-free
remainders. It neither determines $M_n$ for $n\geq4$ nor proves that eight
is the smallest dimension forcing a plane at density $1/10$. Its proof is
independent of polynomial Freiman--Ruzsa theory.
\end{remark}

\section{Entropy first variations and tail exponent rigidity}

All entropies in this note use natural logarithms; the ratio does not depend
on the logarithm base. For independent identically distributed $X,X'$,
write
\[
 \Dplus=H(X+X')-H(X),\quad \Dminus=H(X-X')-H(X),\quad
 \rho=\Dminus/\Dplus.
\]
For every nondegenerate discrete real law of finite entropy,
$\Dplus>0$. Indeed, $\Dplus=I(X;X+X')\geq0$. Equality would imply that the
conditional translated laws of $X+X'$ given any two distinct support points
of $X$ coincide, making the law of $X'$ invariant under a nonzero
translation. The translates of any positive atom would then give infinite
total mass.

Throughout this section the support is \emph{all} of $\N=\{0,1,2,\ldots\}$.
Put $p_n=\Prob(X=n)>0$, and assume
\begin{equation}\label{eq:tail}
 -\log p_n=c n^\beta+o(n^\beta),\qquad c>0,\quad\beta>0.
\end{equation}
This assumption implies finiteness of all polynomial moments and of entropy.
Define the sum and difference laws by
\[
 q_s=\sum_{j=0}^s p_jp_{s-j}\quad(s\geq0),\qquad
 r_m=\sum_{j\geq0}p_jp_{j+|m|}\quad(m\in\mathbb Z).
\]
All these probabilities are positive; $r_{-m}=r_m$.

\begin{lemma}[Coordinate first variation]\label{lem:variation}
For $n\geq1$, let $p^{(t)}=p+t(e_n-e_0)$, with
$|t|<\min(p_0,p_n)$ sufficiently small, where $e_j$ denotes unit mass at $j$.
Define
\begin{equation}\label{eq:gradient}
 G_n=-2\sum_{j\geq0}p_j\log r_{n-j}
       +2\rho\sum_{j\geq0}p_j\log q_{n+j}
       -(\rho-1)\log p_n.
\end{equation}
Then
\begin{equation}\label{eq:derivative}
 \left.\frac{d}{dt}\rho(p^{(t)})\right|_{t=0}
 =\frac{G_n-G_0}{\Dplus(p)}.
\end{equation}
\end{lemma}
\begin{proof}
Holding $\rho=\rho(p)$ fixed while differentiating gives
\[
 \Dplus(p)\,\rho'(0)
 =\left.\frac{d}{dt}\right|_{t=0}
 \bigl[H(r^{(t)})-\rho H(q^{(t)})+(\rho-1)H(p^{(t)})\bigr].
\]
The contribution of varying the $n$th coordinate to the difference entropy
is $-2\sum_jp_j(\log r_{n-j}+1)$, since the two placements of the varied
coordinate give $r_{n-j}$ and $r_{j-n}$, which agree. For the sum entropy
it is $-2\sum_jp_j(\log q_{n+j}+1)$. The input term contributes
$-(\rho-1)(\log p_n+1)$. The coordinate-independent constants cancel
between $n$ and $0$, leaving \eqref{eq:derivative}.

These infinite differentiations are justified as follows. For small enough
$|t|$, $p^{(t)}$ is coordinatewise between $p/2$ and $3p/2$; the output laws
are correspondingly comparable to $q,r$. From \eqref{eq:tail} and the
endpoint terms $p_0p_s$, their negative logarithms are bounded by
$K(1+|s|^\beta)$, uniformly for these $t$. First derivatives of the output
laws are finite linear combinations of translates or reflected translates
of $p$, plus finitely supported terms. They are summable against these
polynomial bounds, since all polynomial moments of $p$ are finite.
Dominated differentiation therefore applies to all three entropy sums.
\end{proof}

\begin{lemma}[Convolution tails]\label{lem:tails}
Under \eqref{eq:tail}, set $\kappa_\beta=\min\{1,2^{1-\beta}\}$. Then
\begin{align}
 \frac{-\log q_n}{c n^\beta}&\longrightarrow\kappa_\beta,
 &\frac{-\log r_n}{c n^\beta}&\longrightarrow1,\label{eq:convtails}\\
 \frac{-\sum_{j\geq0}p_j\log q_{n+j}}{c n^\beta}
 &\longrightarrow\kappa_\beta,
 &\frac{-\sum_{j\geq0}p_j\log r_{n-j}}{c n^\beta}
 &\longrightarrow1.\label{eq:avtails}
\end{align}
\end{lemma}
\begin{proof}
For each $0<\varepsilon<c$, positive constants $a_\varepsilon,b_\varepsilon$
exist such that for every $j\geq0$,
\[
 a_\varepsilon e^{-(c+\varepsilon)j^\beta}
 \leq p_j\leq
 b_\varepsilon e^{-(c-\varepsilon)j^\beta}.
\]
The minimum of $j^\beta+(n-j)^\beta$ over integer $0\leq j\leq n$ is
asymptotic to $\kappa_\beta n^\beta$. For $\beta\leq1$ it occurs at an
endpoint; for $\beta>1$ it occurs at one of the two central integers.
Using one minimizing summand for a lower bound on $q_n$ and all $n+1$
summands for an upper bound gives the first limit in \eqref{eq:convtails},
since $\log(n+1)=o(n^\beta)$ and $\varepsilon$ is arbitrary.

For the difference law,
\[
 p_0p_n\leq r_n=\sum_{j\geq0}p_jp_{n+j}
 \leq b_\varepsilon e^{-(c-\varepsilon)n^\beta}\sum_jp_j.
\]
This proves its claimed tail limit. Also $q_s,r_s\geq p_0p_s$ for $s\geq0$,
so their negative logarithms are bounded by $K(1+s^\beta)$. The integrands
in \eqref{eq:avtails}, after division by $n^\beta$ for $n\geq1$, are
bounded by $K'(1+j^\beta)$, using $|n-j|\leq n+j$ and
$(n+j)^\beta\leq K_\beta(n^\beta+j^\beta)$. This is integrable under $p$.
For each fixed $j$ the required limit follows from \eqref{eq:convtails}.
Dominated convergence gives \eqref{eq:avtails}.
\end{proof}

\begin{theorem}[Tail exponent rigidity and explicit ascent]\label{thm:tail}
For a law satisfying \eqref{eq:tail}, its first-variation coefficients obey
\begin{equation}\label{eq:tailcoefficient}
 \boxed{\ \frac{G_n}{c n^\beta}\longrightarrow
 1+\rho-2\rho\min\{1,2^{1-\beta}\}.\ }
\end{equation}
If this limit is positive, moving sufficiently small positive mass from
zero to any sufficiently distant atom increases $\rho$. If it is negative,
moving sufficiently small positive mass from any sufficiently distant atom
to zero increases $\rho$.

In particular, if $\rho>1$ and the law is stationary under all finite
coordinate mass transfers---as any local maximum under these transfers must
be---then necessarily
\begin{equation}\label{eq:beta}
 \boxed{\ \beta=\log_2\!\left(\frac{4\rho}{1+\rho}\right)>1.\ }
\end{equation}
\end{theorem}
\begin{proof}
Substitute \eqref{eq:avtails} and \eqref{eq:tail} into \eqref{eq:gradient}.
The three coefficients are $2$, $-2\rho\kappa_\beta$, and $\rho-1$;
their sum is the limit in \eqref{eq:tailcoefficient}. Because $G_0$ is fixed
and finite, a nonzero limit determines the eventual sign of $G_n-G_0$.
Equation~\eqref{eq:derivative} gives the two ascent directions. For each
chosen $n$ the transfer can be small enough to preserve positive mass at
both atoms, and it does not change the tail asymptotic.

At a stationary law, $G_n=G_0$ for every $n$. Therefore the limit must
vanish. If $\beta\leq1$ it equals $1-\rho<0$, which is impossible.
For $\beta>1$, solving
$1+\rho=2\rho\,2^{1-\beta}$ gives \eqref{eq:beta}.
\end{proof}

\begin{corollary}[Gaussian tails are never stationary]\label{cor:gaussian}
A full-support law on $\N$ satisfying $-\log p_n\sim c n^2$ has a strict
ratio-increasing finite-coordinate perturbation whenever $\rho>1$.
Its coefficient in \eqref{eq:tailcoefficient} is exactly one.
\end{corollary}
\begin{proof}
For $\beta=2$, $\kappa_\beta=1/2$, so $1+\rho-2\rho\kappa_\beta=1$.
\end{proof}

\begin{remark}
The result concerns \emph{one-sided} full support and the specified
logarithmic tail law. It makes no claim about two-sided Gaussian laws,
finite supports, arbitrary sparse supports, or the existence of an
extremizer having a power-law logarithmic tail. In particular, the
necessary exponent is not a prediction that a global optimizer exists.
\end{remark}

\section{The full sharp constant in prescribed tail classes}

Let $C_*$ be the supremum of $\rho(P)$ over all nondegenerate discrete
real-valued laws with finite entropy. For fixed $\beta,c>0$, let
\[
 \mathcal T_{\beta,c}=
 \{P:p_n>0\ (n\in\N),\ -\log p_n\sim c n^\beta\}.
\]

\begin{theorem}[Universality of each prescribed tail class]\label{thm:universal}
For every $\beta,c>0$,
\begin{equation}\label{eq:sup}
 \boxed{\ C_*=\sup_{P\in\mathcal T_{\beta,c}}\rho(P).\ }
\end{equation}
The same $C_*$ is the supremum over finitely supported nonnegative-integer
laws with rational probabilities.
\end{theorem}
\begin{proof}
We give all approximation steps, since pointwise convergence alone would
not justify entropy convergence.

\emph{Finite truncation.} Given a finite-entropy law $X$, let $F_m$ be an
increasing finite exhaustion of its support and let $X_m$ have the
conditional law given $X\in F_m$. The formula
\[
 H(X_m)=\Prob(X\in F_m)^{-1}
 \sum_{x\in F_m}p_x\log(1/p_x)+\log\Prob(X\in F_m)
\]
shows $H(X_m)\to H(X)$. Let $Z=(X,X')$, and let $E_m$ be the event that both
coordinates lie in $F_m$, of probability $1-\delta_m\to1$. Since $H(Z)<\infty$,
\[
 \delta_m H(Z\mid E_m^c)
 =\sum_{z\in E_m^c}\Prob(Z=z)\log\frac1{\Prob(Z=z)}
   +\delta_m\log\delta_m\longrightarrow0.
\]
For either $W=X+X'$ or $W=X-X'$, entropy decomposition gives
\[
 H(W)=(1-\delta_m)H(W\mid E_m)+\delta_m H(W\mid E_m^c)
       +I(W;\mathbf1_{E_m}).
\]
The second term is bounded by $\delta_m H(Z\mid E_m^c)$, and the last by
the binary entropy $h(\delta_m)$. Thus $H(W\mid E_m)\to H(W)$; the conditional
law is exactly $X_m\pm X'_m$. Positivity of $\Dplus$ gives convergence of
the entropy ratios. This strengthens the elementary truncation principle
used in \cite[Lemma 7.5]{li} to two-sided convergence for the present outputs.

\emph{Integer modelling and rational probabilities.} For a finite real
support $A$, take a basis of its rational linear span, express all its
members in rational coordinates, and clear denominators. This identifies
$A$ additively with a finite subset of $\mathbb Z^r$. An integer linear
form can be chosen injective on the union of this set, its sumset, and
its difference set: finitely many nonzero difference vectors must have
nonzero image, so avoid their finitely many linear hyperplanes. Explicitly,
a form with coefficients $1,M,\ldots,M^{r-1}$ works for any integer $M$
avoiding finitely many polynomial roots. Its image preserves all sum and
difference collisions and hence all three entropies. Translation makes
the support nonnegative. Approximate its positive probabilities by positive
rational probabilities. Entropy is continuous on each relevant finite
simplex, and $\Dplus>0$, so the ratios converge.

\emph{Adding an arbitrarily prescribed tail.} Let $P$ be any finite
nonnegative-integer law. Define a full-support law $Q$ by
\[
 Q(n)=Z^{-1}\exp[-c(n+1)^\beta],\qquad n\geq0.
\]
Its entropy and all moments are finite. For $0<\varepsilon<1$ put
$P_\varepsilon=(1-\varepsilon)P+\varepsilon Q$. Beyond the finite support
of $P$, its logarithmic tail is
\[
 -\log P_\varepsilon(n)=c(n+1)^\beta+\log Z-\log\varepsilon
 =c n^\beta+o(n^\beta).
\]
Hence $P_\varepsilon\in\mathcal T_{\beta,c}$. For a finite mixture of
finite-entropy laws with mixing weights $w_i$, its entropy lies between
$\sum_iw_iH_i$ and $\sum_iw_iH_i+H(w)$. Apply this to the input mixture and
the four-component mixtures describing its sum and difference outputs.
All component entropies are finite, since $H(U\pm V)\leq H(U)+H(V)$.
It follows that all three entropies converge to those of $P$ as
$\varepsilon\downarrow0$, and so do the ratios. Combining the three steps
proves the theorem; the reverse supremum inequality is immediate.
\end{proof}

\begin{corollary}[Same supremum, no Gaussian-tail maximizer]\label{cor:nonattain}
For every $c>0$, the class $\mathcal T_{2,c}$ has supremum $C_*$ but no
member attaining it. More generally, an extremizer in any
$\mathcal T_{\beta,c}$ can exist only when
\begin{equation}\label{eq:betaband}
 \log_2\!\left(\frac{4L}{1+L}\right)<\beta\leq\log_2(8/3),
 \qquad L=1.4051810329.
\end{equation}
The endpoints are approximately $1.224610521648$ and $1.415037499279$.
\end{corollary}
\begin{proof}
The supremum equality is Theorem~\ref{thm:universal}. An attainer would have
$\rho=C_*>L>1$, by \eqref{eq:knownrange}. Finite-coordinate perturbations
remain inside the same class, so Theorem~\ref{thm:tail} forces
$\beta=\log_2(4C_* /(1+C_*))$. This expression is strictly increasing in
$C_*>0$. The known range \eqref{eq:knownrange} gives \eqref{eq:betaband},
which excludes $\beta=2$.
\end{proof}

The supremum theorem does not force any particular tail shape: every
prescribed tail can be added with arbitrarily small mass to an almost
optimal finite law. The stationarity theorem, in contrast, restricts a
single law that actually attains a local maximum. Keeping these two
quantifiers distinct avoids an unjustified existence claim.

\section{Verification record and remaining limitations}

The package contains two standard-library verification programs:
\begin{verbatim}
python verify_midpoint.py
python verify_entropy_variation.py
\end{verbatim}
The first checks all $19{,}683$ labelled two-colour plane partitions,
including the $1{,}371$ admissible ones; $1{,}159$ maximal-direction fibre
and occupied-quotient instances; $640$ plane-free projection instances;
all $296{,}010$ six-point subsets of $\F^3$ for the blocking claim; all
$75{,}816$ spanning five-point subsets for the sumset claim; and all
$17{,}550$ four-point subsets for the stated classification. It also verifies
the sharp sixteen-point construction and the recurrence arithmetic.
Some three-dimensional fibre checks are sampled products, not an
exhaustion of all three-dimensional colourings. The exact five-point check
finds a minimum sumset size of twelve; the written proof uses only eleven.

The entropy program checks $963$ formal first-variation identities using
exact rational coefficients of formal logarithms. It also certifies a
finite fixed-support ascent: on $\{0,\ldots,8\}$, start with integer weights
$2^{64-n^2}$ and transfer probability $10^{-12}$ from zero to six. The ratio
increases from $1.2166978238640320\ldots$ to $1.2166978239078707\ldots$.
This is a finite diagnostic, not a new bound for $C_*$ or a certificate for
an infinite Gaussian-tail law. Rational logarithm enclosures use the
positive atanh series, ninety terms, its geometric tail bound, and outward
rounding. The receipts retain the full weights, intervals, and positive gap.
The infinite-support assertions rest on the proofs of dominated
differentiation and convergence above, not numerical truncation.

The exploratory directory records a time-limited four-dimensional
mixed-integer search that found no feasible point. That outcome is neither
an infeasibility proof nor evidence of an exact extremal value.

The general $O(k\log k)$ midpoint codimension bound and the entropy range
\eqref{eq:knownrange} remain unchanged. The new midpoint theorem does not
settle higher-flat-free remainders or the exact two-colour extremal
function. The tail condition is necessary, not sufficient, and does not
establish an entropy maximizer or rule out irregular or sparse extremizers.

\appendix
\section{Proof of the prior cap classification}\label{app:cap}

A plane cap has at most four points. Through a fixed cap point there are
four lines, each containing at most one further cap point, so initially the
bound is five. Equality would make every line through any cap point meet
the cap twice; no line could meet it once. But three parallel lines would
partition five points into sizes $2,2,1$, a contradiction. For a four-point
cap, put three points at $0,e_1,e_2$. The fourth can only be $e_1+e_2$,
$e_1-e_2$, or $-e_1+e_2$, each giving an affine parallelogram.

No plane section of an admissible two-colour cap remainder can have exactly
three points. Put such points at $0,e_1,e_2$. The coloured points
$e_1+e_2,e_1-e_2,-e_1+e_2$ have these three points as their pairwise
midpoints, forcing three colours.

Now let $W$ be an affine three-space, $s=|R\cap W|$, and let $n_j$ count its
planes with $j$ remainder points. Possible section sizes are $0,1,2,4$.
The $39$ planes, $13$ planes through each point, four planes through a pair,
and one plane through a noncollinear triple give
\begin{align*}
 n_0+n_1+n_2+n_4&=39,& n_1+2n_2+4n_4&=13s,\\
 n_2+6n_4&=4\binom{s}{2},&4n_4&=\binom{s}{3}.
\end{align*}
Three parallel planes show $s\leq12$. Solving gives
\[
 n_4=\tfrac14\binom{s}{3},\quad
 n_1=13s-8\binom{s}{2}+2\binom{s}{3},\quad
 n_0=39-13s+4\binom{s}{2}-\tfrac34\binom{s}{3}.
\]
For $s=3,5,7,11$, $n_4$ is not an integer. For $s=6,8,9$, the value of
$n_1$ is respectively $-2,-8,-3$. For $s=10,12$, the value of $n_0$ is
respectively $-1,-18$. The only possibilities are $s=0,1,2,4$, with
$n_4=1$ in the last case. Thus any four remainder points in a three-space
are coplanar. Four affinely independent remainder points in the original
ambient space would contradict this conclusion in their affine span.
Therefore the whole remainder has affine dimension at most two, and the
plane classification proves Lemma~\ref{lem:cap}. This rederives the
necessary classification from \cite{previous}; no new priority is claimed
for this appendix.

\begin{thebibliography}{9}
\small\setlength{\itemsep}{3pt}
\bibitem{midpoint}
T. Ko\'sciuszko, \emph{Arithmetic Progressions in Midpoint Colourings},
arXiv:2610.11847v1 (2026), Problem 4 and Theorem 3.
\url{https://arxiv.org/html/2610.11847v1}.
\bibitem{li}
R. Li, \emph{The Entropic Sum-Product Phenomenon}, arXiv:2607.29042v1
(2026), Lemma 7.5 and Section 8.
\url{https://arxiv.org/html/2607.29042v1}.
\bibitem{brouwer}
A. E. Brouwer and A. Schrijver, \emph{The blocking number of an affine
space}, Journal of Combinatorial Theory, Series A \textbf{24} (1978),
251--253. \url{https://ir.cwi.nl/pub/9858/9858D.pdf}.
\bibitem{jamison}
R. E. Jamison, \emph{Covering finite fields with cosets of subspaces},
Journal of Combinatorial Theory, Series A \textbf{22} (1977), 253--266.
\url{https://doi.org/10.1016/0097-3165(77)90001-2}.
\bibitem{first}
B2 research continuation I, \emph{Research note}, 13 October 2026.
Supplied artifact \texttt{B2-new-results-2026-10-13.zip}.
\bibitem{previous}
B2 research continuation II, \emph{Two-colour rigidity and boundary
obstructions for entropy extremizers}, 13 October 2026, Theorem 2.1,
Lemma 3.1, and exact entropy certificate. Supplied artifact
\texttt{B2-further-results-2026-10-13.zip}.
\end{thebibliography}
\end{document}
